\documentclass[a4paper,12pt]{article}

% --- Use XeLaTeX font support ---

\usepackage{fontspec}
\setmainfont{TeX Gyre Termes} % nice Times-like font, installed by texlive-core
\setsansfont{TeX Gyre Heros} % optional sans-serif
\setmonofont{Fira Code}       % optional monospaced font
% --- Math packages ---
\usepackage{amsmath, amssymb, amsthm}
\usepackage{mathtools}

% --- Page layout ---
\usepackage[margin=1in]{geometry}

% --- Better lists ---
\usepackage{enumitem}
% --- Hyperlinks for references ---
\usepackage[colorlinks=true, linkcolor=blue, urlcolor=blue]{hyperref}

% Header: fancy layout like the rovided image
\usepackage{fancyhdr}
\setlength{\headheight}{26pt} % increase if fancyhdr warns
\pagestyle{fancy}
\fancyhf{} % clear header/footer
\rhead{\begin{tabular}{r} Mt.: 3052737\end{tabular}}
\chead{\begin{tabular}{c}Lösungen EA 2\end{tabular}}
\lhead{\begin{tabular}{l}61521 Einführung in die \\ Numerische Mathematik SS26 \end{tabular}}
\renewcommand{\headrulewidth}{0.8pt}
\fancyfoot[C]{- \thepage\ -}
% Ensure the plain page style (used by \maketitle) uses the same header

% --- Line spacing ---
\usepackage{setspace} % provides \doublespacing and \onehalfspacing
\usepackage{booktabs}
% --- Optional solutions environment ---
\newenvironment{solution}{
  \par\noindent\textbf{Solution:}\quad
}{\hfill$\blacksquare$\par\vspace{1em}}

% % --- Title info ---
% \title{Aufgabe \#1}
% \author{Your Name}
% \date{\today}
% 
\begin{document}
\raggedright
% Enable double (2.0) line spacing for the document body
\doublespacing
\noindent\textbf{Aufgabe 3} \\
\noindent\textbf{(a)} \\
\begin{center}
\begin{tabular}{ccc}
\toprule
$k$ & klassisch $x_k$ & beschleunigt $\bar{x}_k$ \\
\midrule
0 & 1.200000000000000 & 1.165912859453899 \\
1 & 1.176005207095135 & 1.165596962702548 \\
2 & 1.168783233375686 & 1.165571998094516 \\
3 & 1.166560400718099 & 1.165564582089440 \\
4 & 1.165871558884363 & 1.165562237135375 \\
5 & 1.165657640770319 & 1.165561510564398 \\
\fbox{6} & \fbox{1.165591165577340} & \fbox{1.165561185298521} \\
7 & 1.165570504170872 &   \\
8 & 1.165564081915407 &  \\
\bottomrule
\end{tabular}
\end{center}
Vergleichswert:
$$x^* = 1.165561185207211$$
Fehler bei $k=6$:
$$|x_6 - x^*|\approx 2.9980 \cdot 10^{-5}$$
$$|\overline{x}_6 - x^*|\approx  9.1310 \cdot 10^{-11}$$

\noindent\textbf{(b)}

Setze
$$
e_k := x_k - x^*.
$$
Dann gilt
$$
e_{k+1} = (L + \delta_k) e_k.
$$
Setze
$$
q_k := L + \delta_k.
$$
Dann ist
$$
e_{k+1} = q_k e_k
$$
und wegen $\delta_k \to 0$ gilt
$$
q_k \to L.
$$
Dann ist auch
$$
\Delta x_k = x_{k+1} - x_k = e_{k+1} - e_k = (q_k - 1) e_k.
$$
Weiter:
$$
\Delta^2 x_k = x_{k+2} - 2x_{k+1} + x_k = e_{k+2} - 2e_{k+1} + e_k.
$$
Da
$$
e_{k+2} = q_{k+1} e_{k+1} = q_{k+1} q_k e_k,
$$
folgt
$$
\Delta^2 x_k = (q_{k+1} q_k - 2q_k + 1) e_k.
$$
Damit ist die Aitken-Transformierte
$$
\bar x_k = x_k - \frac{(\Delta x_k)^2}{\Delta^2 x_k}.
$$
Also eingesetzt folgt für den Fehler:
\begin{align*}
\bar x_k - x^* &= e_k - \frac{(q_k - 1)^2 e_k^2}{(q_{k+1} q_k - 2q_k + 1) e_k} \\
 & = e_k \left( 1 - \frac{(q_k - 1)^2}{q_{k+1} q_k - 2q_k + 1} \right).
\end{align*}
Daraus folgt
$$
\frac{\bar x_k - x^*}{x_k - x^*} = 1 - \frac{(q_k - 1)^2}{q_{k+1} q_k - 2q_k + 1}.
$$
Wegen:
$$
q_{k+1} q_k - 2q_k + 1 - (q_k - 1)^2 = q_k (q_{k+1} - q_k).
$$
Kann man das schreiben als:
$$
\frac{\bar x_k - x^*}{x_k - x^*} = \frac{q_k (q_{k+1} - q_k)}{q_{k+1} q_k - 2q_k + 1}.
$$
Jetzt Grenzübergang:
$$
q_k \to L, \qquad q_{k+1} \to L,
$$
also
$$
q_{k+1} - q_k \to 0.
$$
Der Zähler ist dann null:
$$
L \cdot 0 = 0.
$$
Der Nenner geht gegen
$$
L^2 - 2L + 1 = (1 - L)^2.
$$
Da $|L| < 1$, gilt insbesondere $L \neq 1$, also
$$
(1 - L)^2 \neq 0.
$$
Damit folgt
$$
\lim_{k \to \infty} \frac{\bar x_k - x^*}{x_k - x^*} = 0.
$$
Außerdem ist der Nenner $\Delta^2 x_k$ für hinreichend große $k$ nicht null, weil
$$
q_{k+1} q_k - 2q_k + 1 \to (1 - L)^2 \neq 0
$$
und $e_k \neq 0$ vorausgesetzt ist.

Damit existiert die Aitken-transformierte Folge ab einem Index $k_0$, und sie konvergiert tatsächlich schneller.





\end{document}
